% Self-contained manuscript. Compile with pdfLaTeX twice or latexmk -pdf.
\documentclass[11pt]{article}
\usepackage[T1]{fontenc}
\usepackage[utf8]{inputenc}
\usepackage{lmodern}
\usepackage[a4paper,margin=1in]{geometry}
\usepackage{amsmath,amssymb,amsthm,mathtools,bm}
\usepackage{microtype,booktabs,tabularx,array,enumitem}
\usepackage[font=small,labelfont=bf]{caption}
\usepackage{needspace}
\usepackage[round,authoryear]{natbib}
\usepackage{xurl}
\usepackage[hidelinks]{hyperref}
\hypersetup{
  pdftitle={Near-Optimal Proper Bandit Optimization with Smooth Strongly Convex Losses},
  pdfauthor={Colombo Team},
  pdfkeywords={Bandit convex optimization, one-point feedback, strong convexity, log-concave exploration, minimax regret},
  pdfsubject={One-point adversarial bandit convex optimization on general convex domains}}
\numberwithin{equation}{section}
\allowdisplaybreaks[1]
\setlength{\parindent}{1.25em}
\setlength{\parskip}{2pt}
\setlength{\emergencystretch}{1.5em}
\setlength{\textfloatsep}{15pt plus 3pt minus 3pt}
\setlength{\floatsep}{12pt plus 2pt minus 2pt}
\renewcommand{\arraystretch}{1.15}
\setlist[itemize]{leftmargin=1.5em,itemsep=3pt,topsep=4pt}
\setlist[enumerate]{leftmargin=1.8em,itemsep=4pt,topsep=5pt}
\newcommand{\R}{\mathbb R}
\newcommand{\E}{\mathbb E}
\newcommand{\Pp}{\mathbb P}
\newcommand{\Reg}{\operatorname{Reg}}
\newcommand{\Cov}{\operatorname{Cov}}
\newcommand{\tr}{\operatorname{tr}}
\newcommand{\KL}{\operatorname{KL}}
\newcommand{\TV}{\operatorname{TV}}
\newcommand{\vol}{\operatorname{vol}}
\newcommand{\ip}[2]{\langle #1,#2\rangle}
\newcommand{\norm}[1]{\left\|#1\right\|}
\newcommand{\cK}{\mathcal K}
\newcommand{\cF}{\mathcal F}
\newcommand{\cA}{\mathcal A}
\newcommand{\B}{B_2^d}
\newcommand{\diam}{\operatorname{diam}}
\newcommand{\argmin}{\operatorname*{arg\,min}}
\newcounter{algorithm}
\newtheorem{theorem}{Theorem}[section]
\newtheorem{lemma}[theorem]{Lemma}
\newtheorem{proposition}[theorem]{Proposition}
\newtheorem{corollary}[theorem]{Corollary}
\theoremstyle{definition}
\newtheorem{definition}[theorem]{Definition}
\theoremstyle{remark}
\newtheorem{remark}[theorem]{Remark}
\title{Near-Optimal Proper Bandit Optimization\\with Smooth Strongly Convex Losses}
\author{Colombo Team\\[3pt]\texttt{colombo.2026@outlook.com}}
\date{}
\begin{document}
\maketitle
\begin{abstract}
We study one-point adversarial bandit convex optimization with bounded,
$\mu$-strongly convex, $\beta$-smooth losses on a known convex body
$(4d)^{-1}B_2^d\subseteq K\subseteq B_2^d$. Every played action must lie in
$K$, and the learner observes only the loss at that action. Writing
$\kappa=\beta/\mu$, we give an information-theoretic algorithm with expected
static regret $\widetilde O(d\sqrt{\kappa T})$, capped by $T$, uniformly over
comparators on the boundary as well as in the interior. The construction uses
the Gaussian--gauge comparison of Bizeul and Klartag to obtain a feasible
exploration distribution in covariance coordinates. A compactly supported
smoothing density and a separate play density yield a one-point estimator
whose squared covariance-local norm is bounded by $O(d\log^2(ed))$ almost
surely at every round. Combined with quadratic exponential weights, a centered
averaged-tangent inequality preserves the full strong-convexity term and gives
the stated condition-number dependence. We also adapt the classical quadratic
regret lower bound to exact bounded feedback by using compactly supported
random offsets and an adaptive Hellinger argument. This construction applies
even when $\beta=\mu$. For every fixed positive curvature pair, the worst-case
dimension--horizon rate over the nondegenerate admissible domain family is
$\widetilde\Theta(\min\{T,d\sqrt T\})$. All logarithmic factors and parameter
restrictions are explicit. The algorithm uses exact geometric operations;
polynomial-time implementation and matching dependence on a varying condition
number remain open.
\end{abstract}

\noindent\textbf{Keywords.}
Bandit convex optimization; one-point feedback; strong convexity;
log-concave exploration; minimax regret.

\begin{samepage}
\noindent\textbf{AI Usage.}
Nearly the entire research pipeline for this paper was carried out by
Colombo, powered by GPT-5.6 Sol. 
\par
\end{samepage}

\section{Introduction}\label{sec:intro}
In bandit convex optimization, an action has two roles: it incurs a loss and
provides the learner's only information about that loss. Smoothness and strong
convexity help control the cost of exploration, but a constrained domain can
restrict which directions are available. This issue is most apparent near the
boundary, where a symmetric perturbation may leave the feasible set. An
algorithm that queries outside the domain uses a different information
interface; an algorithm that moves the comparator inward must account for the
resulting approximation error. We study the original constrained problem:
every action is feasible, every round supplies one loss value, and the best
fixed comparator may lie anywhere in the domain.

For smooth, strongly convex losses, the barrier method of
\citet{HL14} and the general-domain membership-oracle method of
\citet{Ito20} achieve $\widetilde O(d^{3/2}\sqrt T)$ regret at fixed
curvature parameters. The classical lower bound is
$\Omega(\min\{T,d\sqrt T\})$ \citep{Shamir13}. The smaller dimension
dependence is already attainable on domains with a constant-parameter
self-concordant barrier, such as a Euclidean ball, and under an appropriate
interior-comparator condition \citep{HL14,Ito20}. The remaining issue is
uniformity over general convex bodies without an interior-margin promise.

We obtain this uniform guarantee. For losses in $[0,1]$ that are
$\mu$-strongly convex and $\beta$-smooth, let $\kappa=\beta/\mu$. On the
normalized domain class specified in Section~\ref{sec:model}, our algorithm
satisfies
\begin{equation}\label{eq:intromain}
 \E\Reg_T(x_T^*)
 \leq\min\!\left\{T,\;C d\sqrt{T\left[
 \log^2(ed)\log(edT\kappa)+\kappa\log(eT)\right]}\right\}.
\end{equation}
Here $x_T^*$ is the best fixed action in hindsight, and $C$ is universal.
Together with a bounded quadratic lower bound, this establishes
$\widetilde\Theta_{\mu,\beta}(\min\{T,d\sqrt T\})$ for every fixed
positive curvature pair whenever the admissible domain family is
nondegenerate. The lower bound is attained on a ball; the upper bound is
uniform over the whole domain family.

\paragraph{Contributions.}
The proof has three components.
\begin{enumerate}
 \item \textbf{Feasible exploration with a bounded local estimator.}
 Let $m$ and $\Sigma$ be the mean and covariance of the current log-concave
 state. The Gaussian--gauge comparison of \citet{BK26} identifies a substantial
 Gaussian mass inside a symmetric part of $\Sigma^{-1/2}(K-m)$ at scale
 $O(\log(ed))$. We use this region to construct two densities: one for smoothing,
 and one for play. The play density is proportional to an envelope for the
 smoothing density's gradient, giving a one-point estimator with
 $\widehat g^{\top}\Sigma\widehat g=O(d\log^2(ed))$ pathwise.
 \item \textbf{A regret comparison that preserves curvature.}
 Averaging upper and lower tangent inequalities about the same centered
 smoothing distribution keeps the full strong-convexity term. With quadratic
 exponential weights, this produces \eqref{eq:intromain}. A homothetic
 comparison distribution handles boundary comparators at a logarithmic
 entropy cost.
 \item \textbf{A lower bound in the exact bounded-feedback model.}
 We retain the quadratic regret--information mechanism of \citet{Shamir13}
 while replacing unbounded observation noise by bounded offsets within
 preselected loss functions. An energy-weighted Hellinger argument gives an
 explicit lower bound in curvature and domain radius, including $\beta=\mu$.
 We also characterize the endpoint at which boundedness forces zero regret.
\end{enumerate}

The algorithm builds on the quadratic exponential-weights and covariance
adaptation framework of \citet{Ito20}. The new ingredient is the feasible
pair of densities and its analysis. It replaces a requirement that an entire
exploration ellipsoid remain inside the domain by a distribution whose support
is feasible and whose estimator has a small local norm. The geometric input is
an established external theorem, stated explicitly in
Section~\ref{sec:geometry}.

\paragraph{Scope and organization.}
We prove expected static regret against an oblivious adversary. The known
domain is available for exact integration and sampling, and geometric
computation is uncharged. Thus the result concerns information complexity;
it does not provide a polynomial-time implementation. Section~\ref{sec:model}
states the model and main theorems, and Section~\ref{sec:related} compares them
with prior guarantees. Sections~\ref{sec:geometry}--\ref{sec:algorithm}
construct the exploration rule and online algorithm.
Section~\ref{sec:analysis} proves the upper bound, and
Section~\ref{sec:lower} introduces the lower-bound construction. The appendices
provide the analytic details and complete lower-bound proof.

\section{Model and main results}\label{sec:model}
\subsection{Domain, losses, and feedback}
Let $d,T\geq1$ be integers. All norms are Euclidean, all logarithms are natural,
and $B_2^d$ is the closed unit ball. Consider the domain class
\begin{equation}\label{eq:domain}
 \cK_d=\left\{K\subset\R^d:K\text{ is compact and convex},\quad
                 (4d)^{-1}B_2^d\subseteq K\subseteq B_2^d\right\}.
\end{equation}
The learner knows $K$, $T$, and the regularity constants. The inball fixes
conditioning and supplies a known interior point, but says nothing about the
location of a minimizer.

For $0<\mu\leq\beta$, let $\cF(K;\mu,\beta)$ be the restrictions to $K$ of
functions that are continuously differentiable on a neighborhood of $K$, take
values in $[0,1]$ on $K$, and satisfy, for all $x,y\in K$,
\begin{align}
 f(y)&\geq f(x)+\ip{\nabla f(x)}{y-x}
                   +\frac\mu2\norm{y-x}_2^2,\label{eq:sc}\\
 f(y)&\leq f(x)+\ip{\nabla f(x)}{y-x}
                   +\frac\beta2\norm{y-x}_2^2.\label{eq:sm}
\end{align}
The bounds are required only on $K$. In particular, no globally bounded,
globally strongly convex extension is assumed. Write $\kappa=\beta/\mu$.

An oblivious adversary fixes $f_{1:T}\in\cF(K;\mu,\beta)^T$ before play. At
round $t$, the learner selects $X_t\in K$ based on past observations and its
private randomness, incurs $f_t(X_t)$, and observes exactly that scalar. There
are no gradient observations, additional loss evaluations, or loss queries
outside $K$. Computation on the known domain is uncharged. We call such a
strategy \emph{proper and one-point}.

\subsection{Regret and admissible parameters}
For $u\in K$, define
\[
 \Reg_T(u)=\sum_{t=1}^T\bigl(f_t(X_t)-f_t(u)\bigr),\qquad
 x_T^*=\argmin_{u\in K}\sum_{t=1}^Tf_t(u).
\]
Strong convexity makes $x_T^*$ unique. For a fixed oblivious sequence, it is
also deterministic, a fact used in the boundary-comparator analysis. Let
$\cA(K)$ be the class of measurable proper one-point strategies. The minimax
quantity of interest is taken only over domains with a nonempty loss class.
Set $\cK_d(\mu,\beta)=\{K\in\cK_d:\cF(K;\mu,\beta)\neq\varnothing\}$. Then
\begin{equation}\label{eq:minimax}
 \mathfrak R^*_{T,\mathrm{sc}}(d;\mu,\beta)
 =\sup_{K\in\cK_d(\mu,\beta)}\inf_{\mathcal A\in\cA(K)}
   \sup_{f_{1:T}\in\cF(K;\mu,\beta)^T}
      \E_{\mathcal A}\Reg_T(x_T^*).
\end{equation}
Thus the strategy can depend on the known domain before the loss sequence is
selected.

\begin{proposition}[Compatibility with bounded losses]\label{prop:admissibility}
For any nonempty $\cF(K;\mu,\beta)$,
\begin{equation}\label{eq:diamcurvature}
 \mu\diam(K)^2\leq8,\qquad \mu\leq32d^2.
\end{equation}
The domain family $\cK_d(\mu,\beta)$ is nonempty if and only if
$\mu\leq32d^2$. At $\mu=32d^2$, it contains only
$K=(4d)^{-1}B_2^d$, and the only admissible loss on this body is
$f(x)=\mu\norm{x}_2^2/2$. Consequently the minimax regret is zero at this endpoint.
\end{proposition}
\begin{proof}
Strong convexity at a midpoint and $0\leq f\leq1$ imply
$\mu\norm{x-y}_2^2/8\leq1$ for all $x,y\in K$. Taking the diameter and then
using the inball gives \eqref{eq:diamcurvature}. Conversely, the quadratic
$\mu\norm{x}_2^2/2$ on $(4d)^{-1}B_2^d$ is admissible when $\mu\leq32d^2$.
At equality, put $r_0=(4d)^{-1}$. Applying strong convexity to $\pm r_0v$ for
any unit vector $v$ forces $f(0)=0$ and $f(\pm r_0v)=1$. The first-order
inequalities at zero imply $\nabla f(0)=0$, hence
$f(x)\geq\mu\norm{x}_2^2/2$ throughout $K$. No point of norm exceeding $r_0$
can belong to $K$. On $r_0B_2^d$, the convex function
$f(x)-\mu\norm{x}_2^2/2$ is nonnegative and vanishes on the boundary, so it
vanishes everywhere. Playing zero has zero regret.
\end{proof}

\subsection{Upper and lower bounds}
\begin{theorem}[Uniform proper one-point upper bound]\label{thm:main}
There is a universal constant $C>0$ such that, for every $0<\mu\leq\beta$,
every $d,T\geq1$ with $\mu\leq32d^2$, and $\kappa=\beta/\mu$,
\begin{equation}\label{eq:generalupper}
 \mathfrak R^*_{T,\mathrm{sc}}(d;\mu,\beta)
 \leq\min\!\left\{T,\;C d\sqrt{T\left[
 \log^2(ed)\log(edT\kappa)+\kappa\log(eT)\right]}\right\}.
\end{equation}
In particular this is $\widetilde O(d\sqrt{\kappa T})$, capped by $T$.
The strategy is proper and one-point on every admissible $K\in\cK_d$, and its
bound has no comparator-margin parameter.

More precisely, take the scale $s=16(1+C_{\mathrm g}\log(ed))$ from
Theorem~\ref{thm:gauge}, and set
\begin{align}
 L&=\log\!\bigl(T(8d+2\beta+6\mu+1)\bigr),\qquad h_T=1+\log T,
       \label{eq:generalL}\\
 \Psi&=s^2L+(20\kappa-4)h_T.\label{eq:Psi}
\end{align}
For every fixed comparator $u\in K$, the strategy of
Theorem~\ref{thm:parameter} satisfies the all-horizon bound
\begin{equation}\label{eq:precisegeneral}
 \E\Reg_T(u)\leq\min\!\left\{T,\;5+2d\sqrt{512e\,T\Psi}\right\}.
\end{equation}
No polynomial dependence on a varying parameter is hidden in this formula.
\end{theorem}

\begin{theorem}[Bounded quadratic lower bound]\label{thm:lowergeneral}\label{prop:lower}
Let $0<\mu\leq\beta$ and choose $R\in[(4d)^{-1},1]$ with
\begin{equation}\label{eq:lowerparams}
 M=\mu R^2<2,\qquad \Delta=1-M/2>0,\qquad
 Q=\min\{M,\Delta^2/M\}.
\end{equation}
For every proper one-point strategy on $K=RB_2^d$, there is an admissible
deterministic oblivious loss sequence such that
\begin{equation}\label{eq:lowergeneral}
 \E\Reg_T(x_T^*)\geq\frac{1}{1024\pi}
                  \min\{QT,\;\Delta d\sqrt T\}.
\end{equation}
Each constructed loss has Hessian exactly $\mu I_d$ and values in $[0,1]$ on
$K$, and the common minimizer is interior. Thus the statement includes
$\beta=\mu$. The quantity $\Delta$ is the bounded-loss headroom left by the quadratic
term. An admissible radius exists whenever $\mu<32d^2$.
\end{theorem}

\begin{corollary}[Minimax rate at fixed curvature]\label{cor:fixed}
For every fixed pair $0<\mu\leq\beta<\infty$, there are constants
$c_{\mu}>0$ and $C_{\mu,\beta}<\infty$, independent of $d,T$, such that for
all integers $d,T\geq1$ satisfying $\mu<32d^2$,
\begin{equation}\label{eq:fixedmatching}
 \begin{split}
 c_{\mu}\min\{T,d\sqrt T\}
 &\leq\mathfrak R^*_{T,\mathrm{sc}}(d;\mu,\beta)\\
 &\leq\min\!\left\{T,\;C_{\mu,\beta}d\log(ed)
                                     \sqrt{T\log(edT)}\right\}.
 \end{split}
\end{equation}
Equivalently,
$\mathfrak R^*_{T,\mathrm{sc}}(d;\mu,\beta)
 =\widetilde\Theta_{\mu,\beta}(\min\{T,d\sqrt T\})$
on this nondegenerate domain family. In the simple regime $0<\mu\leq1$,
Theorem~\ref{thm:lowergeneral} with $R=1$ gives, for all $d,T\geq1$,
\begin{equation}\label{eq:unitlower}
 \mathfrak R^*_{T,\mathrm{sc}}(d;\mu,\beta)
 \geq\frac1{4096\pi}\min\{\mu T,d\sqrt T\}.
\end{equation}
\end{corollary}
\paragraph{Interpretation.}
The constants in \eqref{eq:fixedmatching} may depend on the fixed curvature
pair, but not on $d$ or $T$. Appendix~\ref{app:fixed} proves this uniformity
also for large fixed $\mu$. The result matches the dimension and horizon
dependence up to logarithms; a lower bound growing as $\sqrt{\kappa}$ is
not established. Moreover, the lower bound concerns the worst admissible
body, whereas the upper bound applies to every admissible body and every
fixed comparator, including boundary points.


\paragraph{Coordinates and notation.}
For symmetric matrices, $A\preceq B$ denotes the Loewner order; $\tr(A)$ and
$\Cov(V)$ denote trace and covariance. Matrix square roots are the symmetric
positive-definite roots. Affine normalization below is a device for sampling
only: it does not replace the Euclidean assumptions in
\eqref{eq:sc}--\eqref{eq:sm} by assumptions in a transformed norm. We write
$\E_t$ for conditional expectation given the history before round $t$.

\section{Comparison with prior work}\label{sec:related}
\subsection{The closest proper one-point guarantees}
Table~\ref{tab:comparison} keeps the condition number, domain geometry, and
comparator assumptions visible. The rates concern expected static regret;
the prior upper bounds are shown in their stated tuning regimes. Throughout
the table, $D_K=\diam(K)$, $\nu$ is a self-concordant barrier parameter, and
$K'$ denotes a chosen benchmark region.

\begin{table}[!htbp]
\centering
\small
\setlength{\tabcolsep}{4pt}
\caption{Expected regret with proper one-point feedback. Loss magnitudes are
bounded by one and $\kappa=\beta/\mu$. All regret guarantees may be capped
by $T$ in our $[0,1]$ normalization. Conditions below the table specify the
prior horizon and tuning restrictions.}
\label{tab:comparison}
\begin{tabularx}{\textwidth}{@{}>{\raggedright\arraybackslash}p{0.20\textwidth}
 >{\raggedright\arraybackslash}p{0.42\textwidth}
 >{\raggedright\arraybackslash}X@{}}
\toprule
Reference & Regret bound & Geometry and comparator \\
\midrule
\citet[Thm.~10]{HL14}
& $4d\sqrt{(\nu+2\kappa)T\log T}$\newline
  $\quad+2+\beta D_K^2/2$
& Supplied $\nu$-barrier; comparator anywhere in $K$. \\[5pt]
\citet[Cor.~2]{Ito20}
& $O\!\left(d^{3/2}\sqrt{\kappa T\log T}\right)$
& Membership oracle for $K$; comparator anywhere in $K$. \\[5pt]
\citet[Cor.~1]{Ito20}
& $O\!\left(d\sqrt{\kappa T\log T}
           +\dfrac{d^2\log(ed)}{\mu\rho^2}\right)$
& Membership oracle for $K'$; $K'+\rho B_2^d\subseteq K$, $\rho>0$;
  comparator in $K'$. \\[5pt]
Theorem~\ref{thm:main}
& $\widetilde O(d\sqrt{\kappa T})$\newline
  Full logarithms: \eqref{eq:generalupper}.\newline
  Explicit constants: \eqref{eq:precisegeneral}.
& Every admissible $K\in\cK_d$; comparator anywhere in $K$;
  exact geometric operations. \\
\bottomrule
\end{tabularx}
\vspace{6pt}
\begin{minipage}{\textwidth}
\footnotesize
\textit{Tuning and access.}
For Hazan--Levy, $T>1$ and $T/\log T\geq2(\nu+2\kappa)$.
For Ito ($T>1$), use
$\alpha_t=\min\{1/9+\rho\sqrt{t\eta\mu/2},\sqrt d\}$ and require
$\eta\leq\alpha_1/[2d\log(50T)]$.
Corollary~2 takes $\rho=0$ and
$\eta=\sqrt{\kappa\log T/(dT)}$; Corollary~1 takes
$\eta=\sqrt{\kappa\log T/T}$.
Ito provides a membership-oracle implementation. Our theorem holds at every
horizon through the explicit switch in Theorem~\ref{thm:parameter}, without
a running-time guarantee. Hazan--Levy state regularity on all of $\R^d$;
our assumptions impose the curvature inequalities only on $K$.
\end{minipage}
\end{table}

For a general body with a $\nu=O(d)$ barrier, the first row has scale
$\widetilde O(d\sqrt{(d+\kappa)T})$. At fixed curvature, both general-domain
prior bounds therefore have dimension dependence $d^{3/2}$, compared with
$d$ up to logarithms here. On a ball, a constant-parameter barrier already
gives the latter dependence. The improvement is uniform over bodies, and does
not assert a better polynomial rate on every individual domain. The interior
guarantee has the desired leading term but carries an additive cost that
diverges as $\rho\downarrow0$.

Quadratic exponential weights and covariance-based exploration already appear
in \citet{Ito20}. Our feasible sampling construction supplies the missing
local estimator bound without a comparator margin. At varying $\kappa$, the
explicit expressions should be compared directly: in particular, Hazan--Levy
have additive dependence on $\nu$ and $\kappa$ inside the square root. None of
our lower bounds establishes optimal condition-number dependence.

\subsection{Other settings and geometric ingredients}
\paragraph{Barrier parameters in compact rate statements.}
The single-agent analysis of \citet{Ba25} is often summarized as
$\widetilde O(d\sqrt T)$. Their preprint's Appendix~B, Eq.~(B.7), retains
$dL\sqrt T(1+2\nu\log T)$, where $L$ bounds reward magnitudes. A
dimension-dependent barrier parameter must therefore remain visible when
comparing guarantees over general domains. This distinction concerns the
single-agent regret rate; their multi-agent equilibrium results address a
different objective.

\paragraph{Broader loss classes and alternative feedback.}
One-point smoothing is a classical tool in bandit optimization \citep{FKM05}.
For general convex losses, kernel methods give polynomial-time
$\widetilde O(\operatorname{poly}(d)\sqrt T)$ regret \citep{BEL21}, while
\citet{Lat20} gives an information-theoretic
$O(d^{5/2}\sqrt T\log T)$ bound. The online Newton method of
\citet{Fokkema24} obtains an efficient high-probability
$\widetilde O(d^{7/2}\sqrt T)$ adversarial guarantee without positive strong
convexity. The core Bandit Newton Step guarantee of \citet{Suggala24}
permits actions outside the comparator domain. These results concern a broader
loss class, a stronger computational or probabilistic guarantee, or a different
feasibility requirement.

The strongly convex adaptive-regret theorem of \citet[Theorem~1]{Liu25}
assumes that each loss has an unconstrained minimizer in the feasible set.
Their analysis identifies an additional gradient term when that condition is
dropped. The dimension improvements of \citet{Yu26} primarily use two-point
feedback; their one-point extension concerns linear losses on hyperrectangles.
Neither assumption is part of our model.

\paragraph{Lower bounds and convex geometry.}
\citet[Theorem~6]{Shamir13} establishes the quadratic regret--information
tradeoff underlying our lower bound. We adapt it to bounded realized losses
under exact feedback, retaining curvature and radius. This is distinct from
an optimization-error bound for a single quadratic: informative exploration
itself incurs cumulative regret. The newer lower bounds of \citet{RH26}
concern general Lipschitz convex functions with Gaussian observation noise,
rather than our fixed-curvature class of bounded per-round losses.
Finally, the Gaussian--gauge comparison of \citet{BK26} supplies the geometric
input for the upper bound. We use its unconditional $O(\log(ed))$ expectation
comparison for nonsymmetric gauges and general log-concave measures.

\section{A feasible core from Gaussian--gauge comparison}\label{sec:geometry}
This section isolates the geometric input. The aim is to place substantial
Gaussian mass inside a symmetric feasible region after covariance
normalization, without requiring a Euclidean ball of comparable radius to
lie entirely inside that region.

A gauge is a nonnegative convex positively homogeneous function on $\R^d$;
it need not be even. If a convex body $L$ contains the origin in its interior,
its Minkowski gauge is
\[
 \gamma_L(z)=\inf\{a>0:z\in aL\}.
\]
A random vector is isotropic if it has mean zero and covariance $I_d$.
We use the following external result.

\begin{theorem}[Gaussian--gauge comparison, \citealp{BK26}]\label{thm:gauge}
There is a universal $C_{\mathrm g}>0$ such that, for every isotropic
log-concave random vector $Y$ in $\R^d$, every gauge $\varphi$, and
$G\sim N(0,I_d)$,
\begin{equation}\label{eq:gaugecomparison}
 \E\varphi(G)\leq A_d\E\varphi(Y),
 \qquad A_d=C_{\mathrm g}\log(ed).
\end{equation}
\end{theorem}

For $d\geq2$, this is Corollary~1.6 of \citet{BK26}, with its Gaussian
covariance scale converted to a standard-deviation scale by homogeneity.
The nonsymmetric formulation is stated with their Theorem~1.5 and proved in
Theorem~3.5. It applies to general log-concave laws, not only uniform measures.
For $d=1$, apply the two-dimensional statement to $(Y,G')$, with an independent
standard Gaussian $G'$, and to $(x,z)\mapsto\varphi(x)$; the definition permits
this gauge to vanish in some directions. Thus a larger universal constant
covers every dimension. No KLS conjecture is assumed.
\begin{lemma}[A Gaussian-sized symmetric core]\label{lem:coregeometry}
Let $\pi$ be a full-dimensional log-concave probability on $K$, with mean $m$
and covariance $\Sigma\succ0$. Put
\[
 L=\Sigma^{-1/2}(K-m),\qquad
 S=L\cap(-L),\qquad \gamma=\gamma_S=\max\{\gamma_L,\gamma_L(-\cdot)\}.
\]
Fix the deterministic exploration scale
\begin{equation}\label{eq:s}
 s=16(1+A_d).
\end{equation}
Then
\begin{equation}\label{eq:coregeom}
 \E\gamma(G/s)\leq\frac18,
 \qquad \operatorname{Lip}_2(\gamma)\leq s.
\end{equation}
\end{lemma}
\begin{proof}
The mean is in $\operatorname{int}K$: otherwise a supporting hyperplane
through the mean would force a full-dimensional probability onto that
hyperplane. The covariance is positive definite for the same reason. Hence
$L,S$ are compact convex bodies with the origin in their interior.
For $V\sim\pi$, $Y=\Sigma^{-1/2}(V-m)$ is isotropic and log-concave, and
$\gamma_L(Y)\leq1$ almost surely. Thus Theorem~\ref{thm:gauge} gives
$\E\gamma_L(G)\leq A_d$. Gaussian symmetry implies
\[
 \E\gamma(G)
 \leq\E\gamma_L(G)+\E\gamma_L(-G)\leq2A_d.
\]
This proves the first assertion. Since $\gamma$ is a norm, it is the support
function of the symmetric polar $S^\circ$. For each $v\in S^\circ$,
$\gamma(G)\geq|\ip vG|$, whence
\[
 \operatorname{Lip}_2(\gamma)
 =\sup_{v\in S^\circ}\norm v_2
 \leq\sqrt{\pi/2}\,\E\gamma(G)\leq s.
\]
No Euclidean regularity assumption on the loss was changed in this argument.
\end{proof}

\section{Score-dominating proper exploration}\label{sec:sampler}
A smoothing density supplies a gradient identity, while a play density
determines which action is observed. We choose the play density proportional
to an envelope for the gradient of the smoothing density.
This controls the importance weight even near the boundary and gives a
pathwise bound on the estimator.

For the objects in Lemma~\ref{lem:coregeometry}, let $\phi_s$ be the density of
$N(0,s^{-2}I_d)$, and define
\begin{align}
 h(z)&=(1-\gamma(z))_+,\label{eq:h}\\
 Z&=\int_{\R^d}\phi_s(z)h(z)^2\,dz,\qquad
 p(z)=\frac{\phi_s(z)h(z)^2}{Z},\label{eq:p}\\
 N&=\int_{\R^d}\phi_s(z)h(z)(s^2\norm z_2+2s)\,dz,\qquad
 q(z)=\frac{\phi_s(z)h(z)(s^2\norm z_2+2s)}{N}.\label{eq:q}
\end{align}
Here $p$ is a smoothing density and $q$ is the play density. Both are determined by the known domain and the current mean and covariance.
The taper $h$ vanishes at the boundary; its square in $p$ makes the
integration-by-parts identity compatible with a nonsmooth support.
Draw $z\sim q$, play $X=m+\Sigma^{1/2}z$, and, after observing $f(X)$, set
\begin{equation}\label{eq:estimator}
 \widehat g=\frac NZ f(X)\Sigma^{-1/2}
 \frac{s^2 z h(z)+2\nabla\gamma(z)}{s^2\norm z_2+2s}.
\end{equation}
The gradient of the gauge exists Lebesgue almost everywhere. We set $\nabla\gamma=0$ on its
null nondifferentiability set. This is a
\emph{domain-geometric} derivative, not a loss-gradient oracle.

\begin{lemma}[Sampler guarantees]\label{lem:sampler}
The above sampler has all the following properties:
\begin{align}
 &X\in K\quad\text{almost surely},\qquad \E_p[z]=\E_q[z]=0;\label{eq:feasible}\\
 &\E_p[zz^\top]\preceq\frac8{s^2}I_d,
 \qquad \E_q[zz^\top]\preceq\frac{32}{s^2}I_d;\label{eq:explorecov}\\
 &\E_q[\widehat g]
     =\E_{z\sim p}[\nabla f(m+\Sigma^{1/2}z)];\label{eq:unbiased}\\
 &\norm{\Sigma^{1/2}\widehat g}_2\leq16s\sqrt d
   \quad\text{almost surely, whenever }0\leq f\leq1.\label{eq:boundedlocal}
\end{align}
Every term used by the sampler can be specified using domain geometry and one
loss evaluation at the played point.
\end{lemma}
\begin{proof}
All densities in \eqref{eq:p}--\eqref{eq:q} vanish outside $S$, and
$m+\Sigma^{1/2}S\subseteq K$. They are even, so both centered perturbations
have mean zero. This proves \eqref{eq:feasible}.

Let $G\sim N(0,I_d)$. Markov's inequality and \eqref{eq:coregeom} give
\[
 \Pp(\gamma(G/s)\leq1/2)\geq3/4,
 \qquad Z\geq\frac{3}{16}.
\]
Also
\[
 N=s\E\bigl[h(G/s)(\norm G_2+2)\bigr]\leq3s\sqrt d.
\]
For a lower bound, denote the event $\{\gamma(G/s)>1/2\}$ by $B$. Then
\[
 \E[\norm G_2\mathbf1_B]\leq\sqrt{d\Pp(B)}\leq\tfrac12\sqrt d,
 \qquad \E\norm G_2\geq\sqrt{2/\pi}\sqrt d.
\]
The latter follows from $\norm G_2\geq d^{-1/2}\sum_i|G_i|$. Consequently,
\begin{equation}\label{eq:Nlower}
 N\geq b\,s\sqrt d,
 \qquad b=\tfrac12\bigl(\sqrt{2/\pi}-\tfrac12\bigr)>0.
\end{equation}
Since $h\leq1$ and $Z\geq3/16$, the $p$ covariance bound is immediate.
For any $A\succeq0$, Gaussian moments imply
\[
 \E[(G^\top A G)^2]\leq3(\tr A)^2,
 \qquad \E[(G^\top A G)\norm G_2]\leq\sqrt{3d}\,\tr A.
\]
Thus
\[
 \E_q[z^\top A z]
 \leq\frac{(\sqrt{3d}+2)\tr A}{sN}
 \leq\frac{\sqrt3+2}{b}\frac{\tr A}{s^2}
 \leq\frac{32\tr A}{s^2},
\]
which proves \eqref{eq:explorecov}.

The gauge is globally Lipschitz. Therefore $h=(1-\gamma)_+$ is
Lipschitz and compactly supported, and $p\in W^{1,1}(\R^d)$ (indeed it is
Lipschitz with compact support). This is a \emph{global} zero-extension statement,
not an appeal to a smooth boundary of $S$. In particular, its distributional
gradient is an $L^1$ vector field and has no singular boundary measure.
Almost everywhere in $S$,
\begin{equation}\label{eq:scoreidentity}
 -\nabla p(z)=\frac{\phi_s(z)}{Z}
       \bigl(s^2z h(z)^2+2h(z)\nabla\gamma(z)\bigr).
\end{equation}
For $q$-almost every $z$, the estimator \eqref{eq:estimator} is exactly
\[
 f(m+\Sigma^{1/2}z)\Sigma^{-1/2}\frac{-\nabla p(z)}{q(z)}.
\]
Integration by parts therefore gives
\[
 \int q(z)\widehat g(z)\,dz
 =\Sigma^{-1/2}\int p(z)\nabla_z f(m+\Sigma^{1/2}z)\,dz
 =\int p(z)\nabla f(m+\Sigma^{1/2}z)\,dz.
\]
For complete justification, multiply the $C^1$ function
$z\mapsto f(m+\Sigma^{1/2}z)$ on a neighborhood of $S$ by a smooth compactly
supported cutoff equal to one near $S$, and extend by zero. Apply the weak
integration-by-parts identity to this test function. Only its values and
first derivatives on $S$ survive. No loss evaluation outside $K$ is used.
All of $m,\Sigma,p,q$ are held fixed in this identity: we are \emph{not}
differentiating a family of distributions depending on $m$.

Finally, $\norm{\nabla\gamma(z)}_2\leq s$ almost everywhere, and $h\leq1$.
The vector fraction in \eqref{eq:estimator} has Euclidean norm at most one.
Since $N/Z\leq16s\sqrt d$, \eqref{eq:boundedlocal} follows.
\end{proof}

\begin{remark}[Sampling before evaluating the loss]
The play law is a genuine probability kernel. For example, propose from the
normalized density proportional to
$\phi_s(z)(s^2\norm z_2+2s)$, then accept with probability $h(z)$.
Its acceptance probability is $N/[s(\E\norm G_2+2)]\geq b/3>0$.
The radial proposal is a mixture of a Gaussian radius and its radius-biased
law. Rejected proposals require geometry only; the single loss evaluation is
made \emph{after} acceptance. The ratio $N/Z$ and the mean and covariance
remain exact geometric integrals here. This observation does not establish
an efficient finite-precision implementation of the complete algorithm.
\end{remark}


\paragraph{The role of the play density.}
The score $-\nabla\log p$ contains a term proportional to
$\nabla\gamma/h$, which can diverge as $h$ approaches zero. Direct sampling
from $p$ would not provide \eqref{eq:boundedlocal}. In contrast, $q$ has only
one power of $h$ and includes a radial upper bound on the remaining score.
Consequently $(-\nabla p)/q$ has bounded Euclidean norm. The construction does
not identify $p$ and $q$: the smoothing bias below uses $p$, while the cost of
the action actually played uses $q$. Exact support in $S$ is also stronger
than a high-probability feasibility argument.

\section{The online algorithm}\label{sec:algorithm}
Set $\lambda=\mu$ and fix a learning rate $\eta>0$. We maintain quadratic
potentials and their exponential-weights distributions:
\begin{align}
 F_1(x)&=\frac\lambda2\norm{x}_2^2,\label{eq:F1}\\
 \pi_t(dx)&=\frac{e^{-\eta F_t(x)}\,dx}{\int_K e^{-\eta F_t(v)}\,dv},
 \qquad m_t=\E_{\pi_t}x,\quad\Sigma_t=\Cov_{\pi_t}(x),\label{eq:pi}\\
 \widehat\ell_t(x)&=\ip{\widehat g_t}{x-m_t}
                  +\frac\lambda2\norm{x-m_t}_2^2,
 \qquad F_{t+1}=F_t+\widehat\ell_t.\label{eq:update}
\end{align}
Only a single tilt vector needs to be updated: up to an irrelevant
additive constant,
\[
 F_t(x)=\frac{\lambda t}{2}\norm{x}_2^2+\ip{b_t}x,
 \qquad b_{t+1}=b_t+\widehat g_t-\lambda m_t.
\]
Algorithm~\ref{alg:main} specifies the complete interaction protocol. Its inputs
include $s$ from \eqref{eq:s}. The learning-rate selection and the option of
using a trivial strategy at small horizons are specified in
Theorem~\ref{thm:parameter} below.

\begin{figure}[tbp]
\begin{minipage}{\textwidth}
\small
\hrule\vspace{6pt}
\refstepcounter{algorithm}\label{alg:main}
\noindent\textbf{Algorithm \thealgorithm. Proper Gaussian-core exponential weights}
\par\vspace{5pt}\hrule\vspace{6pt}
\noindent\textbf{Input:} Known $K\in\cK_d$, horizon $T$, $\lambda=\mu$,
scale $s$, learning rate $\eta$.\par
\noindent\textbf{Initialize:} $b_1=0$.\par
\noindent\textbf{For each round $t=1,\ldots,T$:}
\begin{enumerate}[label=\arabic*.,leftmargin=1.8em,itemsep=3pt,topsep=4pt]
 \item Form $\pi_t(dx)\propto
   \exp\!\left(-\eta[\lambda t\norm{x}_2^2/2+\ip{b_t}x]\right)\mathbf1_K(x)\,dx$.
 \item Compute its mean $m_t$ and covariance $\Sigma_t$.
 \item Set $L_t=\Sigma_t^{-1/2}(K-m_t)$, $S_t=L_t\cap(-L_t)$,
   and $\gamma_t=\gamma_{S_t}$.
 \item Define $h_t,Z_t,N_t,p_t,q_t$ by \eqref{eq:h}--\eqref{eq:q}.
 \item Draw $z_t\sim q_t$ and play $X_t=m_t+\Sigma_t^{1/2}z_t\in K$.
 \item Observe the single scalar $y_t=f_t(X_t)$.
 \item Form $\widehat g_t$ by \eqref{eq:estimator}, using $y_t$.
 \item Set $b_{t+1}=b_t+\widehat g_t-\lambda m_t$.
\end{enumerate}
\vspace{2pt}\hrule
\end{minipage}
\end{figure}

The history $\mathcal H_{t-1}$ contains all previous actions, observations, and
internal state. For a fixed oblivious loss sequence, $f_t$ is deterministic,
and $\pi_t,m_t,\Sigma_t,p_t,q_t$ are $\mathcal H_{t-1}$-measurable. Conditional
on this history, Lemma~\ref{lem:sampler} applies with no random parameters left
unconditioned. Its integrals, rejected proposals, and gauge derivatives are
geometric operations, not additional observations of $f_t$.

The next two lemmas connect this state to the regret proof.
The proofs, including finite-horizon integrability, appear in
Appendix~\ref{app:analytic}.

\begin{lemma}[Covariance contraction and well-definedness]\label{lem:state}
Algorithm~\ref{alg:main} defines a measurable proper one-point strategy.
For each fixed finite horizon its estimators and surrogates are integrable.
Moreover,
\begin{equation}\label{eq:covcontraction}
 \Sigma_t\preceq\frac1{\eta\lambda t}I_d,
 \qquad \sum_{t=1}^T\tr\Sigma_t
       \leq\frac{d(1+\log T)}{\eta\lambda}.
\end{equation}
These statements hold on arbitrary convex bodies, including nonsmooth ones.
\end{lemma}

\begin{lemma}[Local exponential moments]\label{lem:mgf}
For a log-concave probability $\pi$ with mean $m$ and covariance $\Sigma$, if
$v^\top\Sigma v\leq1/4$, then
\begin{equation}\label{eq:mgfexplicit}
 \log\E_{U\sim\pi} e^{\ip v{U-m}}\leq2e\,v^\top\Sigma v.
\end{equation}
\end{lemma}

This local inequality is sufficient even when $\widehat g_t$ has a large
ordinary Euclidean norm. By \eqref{eq:boundedlocal}, every realized estimator
satisfies the hypothesis for $v=-\eta\widehat g_t$ as soon as
\begin{equation}\label{eq:etarestriction}
 \eta\leq\frac1{32s\sqrt d}.
\end{equation}

\section{Regret analysis}\label{sec:analysis}
The proof has two separate steps: an exponential-weights bound for the quadratic
surrogates, including boundary comparators, and a smoothness-based comparison
between those surrogates and the actual incurred losses.

\subsection{Surrogate regret with boundary comparators}
A simple consequence of the loss assumptions is
\begin{equation}\label{eq:globalG}
 \sup_{x\in K}\norm{\nabla f_t(x)}_2\leq G_0:=4d+\beta.
\end{equation}
Indeed, convexity and the inball bound $\norm{\nabla f_t(0)}_2$ by $4d$, and
smoothness controls the variation away from zero. Appendix~\ref{app:regularity}
justifies the gradient-Lipschitz consequence of \eqref{eq:sm} on the closed
domain. This gradient bound is used only in the comparison-distribution error;
its dimension dependence enters the regret bound through a logarithm.
\begin{lemma}[Expected surrogate bound including boundary comparators]
\label{lem:surreg}
Put $G_0=4d+\beta$ and
\begin{equation}\label{eq:delta}
 \delta=\frac1{T(2G_0+6\lambda+1)}.
\end{equation}
If $\eta\leq1/(32s\sqrt d)$, then, for every fixed $u\in K$,
\begin{equation}\label{eq:surreg}
 \E\left[-\sum_{t=1}^T\widehat\ell_t(u)\right]
 \leq5+\frac{d\log(1/\delta)}\eta+C_1\eta d s^2T,
\end{equation}
where one may take $C_1=512e$.
\end{lemma}
\begin{proof}
Let $U_K$ be uniform on $K$, and let
$\mathcal Z_t=\int e^{-\eta F_t(x)}U_K(dx)$. For every realized estimator,
\begin{align*}
 \frac{\mathcal Z_{t+1}}{\mathcal Z_t}
 &=\E_{V\sim\pi_t}
     \exp\left(-\eta\ip{\widehat g_t}{V-m_t}
                     -\frac{\eta\lambda}2\norm{V-m_t}_2^2\right)\\
 &\leq\E_{V\sim\pi_t}e^{-\eta\ip{\widehat g_t}{V-m_t}}.
\end{align*}
Lemma~\ref{lem:sampler} and Lemma~\ref{lem:mgf} imply
\[
 \log\mathcal Z_{T+1}-\log\mathcal Z_1
 \leq 2e\eta^2\sum_{t=1}^T\widehat g_t^\top\Sigma_t\widehat g_t
 \leq512e\eta^2ds^2T.
\]
The integration variable $V$ is an auxiliary draw used only to evaluate the
partition-function ratio, not an extra loss query. This step is pathwise:
the exponential-moment inequality holds for every
vector of the prescribed local norm, even though $\widehat g_t$ is random.
Since $F_1\geq0$, $\log\mathcal Z_1\leq0$.

For any deterministic comparison probability $Q\ll U_K$, the variational
inequality for the log partition function gives
\[
 \log\mathcal Z_{T+1}\geq-\eta\E_Q F_{T+1}-\KL(Q\Vert U_K).
\]
Therefore,
\begin{equation}\label{eq:Qbound}
 -\sum_{t=1}^T\E_Q\widehat\ell_t
 \leq\E_QF_1+\frac{\KL(Q\Vert U_K)}\eta+C_1\eta ds^2T.
\end{equation}
Use $Q=Q_{u,\delta}$, the uniform law on $(1-\delta)u+\delta K$. This set is
contained in $K$ even for $u\in\partial K$, and
\[
 \KL(Q_{u,\delta}\Vert U_K)=d\log(1/\delta).
\]
All $x$ in its support satisfy $\norm{x-u}_2\leq2\delta$. Since $m_t,u\in K$,
\[
 \left|\frac\lambda2
 (\norm{x-m_t}_2^2-\norm{u-m_t}_2^2)\right|\leq6\lambda\delta.
\]
By \eqref{eq:globalG}, the loss gradients satisfy
$\norm{\nabla f_t(x)}_2\leq G_0$ on $K$. By \eqref{eq:unbiased},
$\norm{\E_t\widehat g_t}_2\leq G_0$. Consequently,
\[
 \E\left[\E_Q\widehat\ell_t(x)-\widehat\ell_t(u)\right]
 \leq(2G_0+6\lambda)\delta.
\]
The comparison law $Q_{u,\delta}$ is deterministic. Lemma~\ref{lem:state}
justifies Fubini and the tower property in this calculation. Sum the inequality,
take expectations in \eqref{eq:Qbound}, and use
$\E_QF_1\leq4$ and \eqref{eq:delta}. Indeed, $\lambda=\mu$, $0\in K$,
and Proposition~\ref{prop:admissibility} imply
$F_1(x)=\mu\norm{x}_2^2/2\leq\mu\diam(K)^2/2\leq4$.
This proves
\eqref{eq:surreg}. Comparator smoothing is used only in this analysis;
it does not add a loss-oracle call or impose a margin on $u$.
\end{proof}


\subsection{An averaged-tangent comparison}
The following centered comparison avoids a norm bound on the gradient bias.
It retains the full strong-convexity term, so we may use
$\lambda=\mu$ rather than sacrifice part of the curvature in Young's inequality.

\begin{lemma}[Centered averaged tangents]\label{lem:tangent}
Let $f$ be $\mu$-strongly convex and $\beta$-smooth on $K$. Let $Y$ be a
$K$-valued random vector with $\E Y=m\in K$, write
$S_Y=\E\norm{Y-m}_2^2$ and $g_Y=\E\nabla f(Y)$. Then for every $u\in K$,
\begin{equation}\label{eq:tangent}
 f(m)-f(u)\leq\ip{g_Y}{m-u}-\frac\mu2\norm{m-u}_2^2
                         +\frac{\beta-\mu}{2}S_Y.
\end{equation}
\end{lemma}
\begin{proof}
Apply the smooth upper inequality from $Y$ to $m$ and the strongly convex
lower inequality from $Y$ to $u$, then subtract:
\[
 f(m)-f(u)\leq\ip{\nabla f(Y)}{m-u}
            +\frac\beta2\norm{m-Y}_2^2
            -\frac\mu2\norm{u-Y}_2^2.
\]
Take expectations and use
$\E\norm{u-Y}_2^2=\norm{u-m}_2^2+S_Y$. All points in these inequalities
belong to $K$, even when $u\in\partial K$.
\end{proof}

\begin{lemma}[True regret from the two densities]\label{lem:regtransfer}
For every fixed $u\in K$ and every round, Algorithm~\ref{alg:main} with
$\lambda=\mu$ satisfies
\begin{equation}\label{eq:regtransfer}
 \E_t[f_t(X_t)-f_t(u)]
 \leq-\E_t\widehat\ell_t(u)
          +(20\beta-4\mu)\frac{\tr\Sigma_t}{s^2}.
\end{equation}
\end{lemma}
\begin{proof}
Condition on the past. Write $m=m_t$, $\Sigma=\Sigma_t$,
$\bar g=\E_t\widehat g_t$, and let $Y=m+\Sigma^{1/2}z$ with $z\sim p$ as an
auxiliary integration variable, not a played action. Define
\[
 S_p=\E_p\norm{\Sigma^{1/2}z}_2^2,\qquad
 S_q=\E_q\norm{\Sigma^{1/2}z}_2^2.
\]
The sampler guarantees $\E Y=\E_tX_t=m$,
$\bar g=\E_p\nabla f_t(Y)$, and
$S_p\leq8\tr\Sigma/s^2$, $S_q\leq32\tr\Sigma/s^2$.
By smoothness of the actual play and Lemma~\ref{lem:tangent},
\begin{align*}
 \E_t[f_t(X_t)-f_t(u)]
 &\leq f_t(m)-f_t(u)+\frac\beta2S_q\\
 &\leq\ip{\bar g}{m-u}-\frac\mu2\norm{m-u}_2^2
                +\frac{\beta-\mu}{2}S_p+\frac\beta2S_q\\
 &=-\E_t\widehat\ell_t(u)
                +\frac{\beta-\mu}{2}S_p+\frac\beta2S_q.
\end{align*}
Substitution of the two second-moment bounds yields
$4(\beta-\mu)+16\beta=20\beta-4\mu$.
\end{proof}

A direct bound on $\norm{\bar g-\nabla f_t(m)}$ followed by Young's inequality
would instead introduce a $\beta^2/\mu$ term before covariance contraction.
Lemma~\ref{lem:tangent} avoids that unnecessary loss. The two densities remain
distinct: $p$ supplies the averaged tangent and $q$ supplies the incurred cost.
No translated distribution around the comparator is required.

\Needspace{16\baselineskip}
\subsection{Parameter-explicit guarantee and horizon tuning}
\begin{theorem}[Explicit upper bound and legal tuning]\label{thm:parameter}
Fix $0<\mu\leq\beta$, $K\in\cK_d(\mu,\beta)$, and an oblivious admissible
loss sequence. Set $\lambda=\mu$, $\kappa=\beta/\mu$, $G_0=4d+\beta$, and
\begin{align}
 \delta&=\frac1{T(2G_0+6\mu+1)},\nonumber\\
 H&=\log(1/\delta)+\frac{(20\kappa-4)(1+\log T)}{s^2},\label{eq:H}\\
 U&=5+2ds\sqrt{512e\,HT}.\label{eq:U}
\end{align}
If $U<T$, run Algorithm~\ref{alg:main} with
$\eta=\sqrt{H/(512e\,s^2T)}$; otherwise play $0$ at every round. Then, for
every deterministic $u\in K$ and every horizon,
\begin{equation}\label{eq:explicitupper}
 \E\Reg_T(u)\leq\min\{T,U\}.
\end{equation}
\end{theorem}
\begin{proof}
Combine Lemmas~\ref{lem:surreg} and~\ref{lem:regtransfer} with covariance
contraction \eqref{eq:covcontraction}. Since $\lambda=\mu$, every
$\eta\leq1/(32s\sqrt d)$ gives the master inequality
\begin{equation}\label{eq:master}
 \E\Reg_T(u)\leq5+
 \frac d\eta\left[\log(1/\delta)
                  +\frac{(20\kappa-4)(1+\log T)}{s^2}\right]
                 +512e\,\eta ds^2T.
\end{equation}
In particular, the curvature coefficient is exactly
$(20\beta-4\mu)/\mu=20\kappa-4$, with all traces measured in the original
Euclidean coordinates.

If $U<T$, then $T>4(512e)d^2s^2H$. This implies $T\geq(2/e)dH$, which is
precisely sufficient for
$\eta=\sqrt{H/(512e\,s^2T)}\leq1/(32s\sqrt d)$.
Substitution in \eqref{eq:master} gives $\E\Reg_T(u)\leq U$.
Otherwise the fixed action zero gives regret at most $T$. Thus no hidden
large-horizon condition is needed.
\end{proof}

\begin{proof}[Proof of Theorem~\ref{thm:main}]
The identity $s^2H=\Psi$ makes \eqref{eq:precisegeneral} exactly the conclusion
of Theorem~\ref{thm:parameter}. To derive \eqref{eq:generalupper}, note that
Proposition~\ref{prop:admissibility} gives $\mu\leq32d^2$, and therefore
$\beta=\kappa\mu\leq32\kappa d^2$. Consequently
\[
 8d+2\beta+6\mu+1\leq265\kappa d^2,
 \qquad L\leq C\log(edT\kappa).
\]
Since $s=O(\log(ed))$, $20\kappa-4\leq20\kappa$, and $d,T,\kappa\geq1$,
a universal constant absorbs the additive five in \eqref{eq:precisegeneral}.
This proves \eqref{eq:generalupper} for each fixed $u\in K$. Set $u=x_T^*$,
which is deterministic for an oblivious sequence, and take the minimax extrema.
\end{proof}

\paragraph{Where the dimension improvement occurs.}
The partition-function argument pays an entropy cost of
$O(d\log(edT))/\eta$ and an estimation cost of $O(\eta ds^2T)$.
The exploration error contributes
$O(d\kappa(1+\log T)/(\eta s^2))$. At fixed curvature, the balance is $d s\sqrt{T\log(edT)}$.
The only nonconstant geometric scale is $s=O(\log(ed))$; no barrier parameter,
comparator margin, or inverse covariance eigenvalue enters this final bound.
The finite-horizon eigenvalue bounds needed for integrability are deliberately
kept out of the rate calculation.

\section{Quadratic lower bound under exact bounded feedback}\label{sec:lower}
We use quadratic losses to retain the curvature parameters without imposing a
positive gap $\beta-\mu$. The argument adapts the regret--information tradeoff of
\citet[Theorem~6]{Shamir13} to exact bounded feedback. Compactly supported
offsets keep every realized loss admissible. Since shifted compact supports
need not be mutually absolutely continuous, we control Hellinger distance
instead of Kullback--Leibler divergence.

Fix the radius and bounded-loss headroom in \eqref{eq:lowerparams}, and put
\begin{equation}\label{eq:quadscale}
 \sigma=\Delta/8,\qquad
 a^2=\min\!\left\{\frac{\mu^2R^2}{64d},\;
                   \frac{\Delta^2}{64dR^2},\;
                   \frac{\mu\sigma}{16\pi\sqrt T}\right\}.
\end{equation}
Let $\theta$ be uniform on $\{-1,1\}^d$, and presample independent offsets
$\xi_t$ with density
$q_0(z)=\cos^2(\pi z/2)\mathbf1_{\{|z|\leq1\}}$. Define
\begin{equation}\label{eq:lowerfamily}
 f_t(x)=b_0+\frac\mu2\norm{x}_2^2-a\ip{\theta}x+\sigma\xi_t,
 \qquad b_0=\sigma+aR\sqrt d.
\end{equation}
The Hessian is exactly $\mu I_d$, so these losses belong to every smoothness
class with $\beta\geq\mu$. The scale in \eqref{eq:quadscale} keeps all values
inside $[0,1]$ and places the common minimizer
$x_\theta^*=(a/\mu)\theta$ inside $(R/8)B_2^d$.

A single sign flip changes the observation location by $2a x_i$. In contrast
to a nonlinear testing family with uniformly small value differences, this
difference can be larger at informative exploration points. The correct
argument therefore tracks \emph{query energy}: the average Hellinger distance
is controlled by $a^2\sum_t\E\norm{X_t}_2^2$, while the quadratic regret
controls that same energy. A learner cannot acquire the signs cheaply by
moving away from the minimizer. Appendix~\ref{app:lower} supplies the complete
weighted recursion and contradiction argument, yielding
\eqref{eq:lowergeneral}. The offsets cancel from regret and are part of a
sequence fixed before play, not an additional noise channel.

\section{Discussion}\label{sec:discussion}
The fixed-curvature dimension--horizon rate of proper one-point bandit
optimization is $\widetilde\Theta(\min\{T,d\sqrt T\})$ on the nondegenerate
admissible domain family. The upper bound is uniform at the boundary and
requires no dimension-dependent barrier parameter. Its main geometric feature
is a feasible distribution with substantial mass, rather than a fully contained
exploration ellipsoid. Separating the smoothing and play densities converts
this feature into a pathwise local estimator bound.

Three questions remain. First, the algorithm uses exact moments, gauge
geometry, and normalizing integrals. The constant-acceptance sampler handles
only the final draw; an efficient implementation would also need to approximate
the state and its geometry while controlling feasibility and estimation error.
Second, the upper and lower bounds leave logarithmic factors and the dependence
on varying curvature parameters unmatched. The lower construction includes
$\beta=\mu$, but does not yield a growing condition-number factor. The
zero-regret endpoint further rules out a positive lower constant uniform over
all feasible curvature choices.

Third, the theorem is an expected-regret statement for an oblivious loss
sequence. The boundary-comparator step controls the conditional mean of the
estimator, rather than a uniformly small Euclidean magnitude of every
realization. High-probability regret and a fully adaptive adversary require
additional arguments. The present proof establishes the information-theoretic
rate under the stated protocol, with these computational and probabilistic
extensions left open.

\clearpage
\begin{thebibliography}{14}
\small
\setlength{\itemsep}{5pt}
\raggedright

\bibitem[Ba et~al.(2025)]{Ba25}
Wenjia Ba, Tianyi Lin, Jiawei Zhang, and Zhengyuan Zhou.
\newblock Doubly optimal no-regret online learning in strongly monotone games
with bandit feedback.
\newblock \emph{Operations Research}, 73(6):3219--3244, 2025.
\newblock doi:10.1287/opre.2021.0445.
\newblock Parameter-explicit comparison uses Appendix~B, Eq.~(B.7), of the
accompanying preprint, \url{https://arxiv.org/abs/2112.02856v4}.

\bibitem[Bizeul and Klartag(2026)]{BK26}
Pierre Bizeul and Boaz Klartag.
\newblock Distances between non-symmetric convex bodies: optimal bounds up to
polylog.
\newblock arXiv:2510.20511v3, July 1, 2026. First circulated in 2025.
\newblock \url{https://arxiv.org/abs/2510.20511v3}.

\bibitem[Bubeck et~al.(2021)]{BEL21}
S\'ebastien Bubeck, Ronen Eldan, and Yin Tat Lee.
\newblock Kernel-based methods for bandit convex optimization.
\newblock \emph{Journal of the ACM}, 68(4), Article 25:1--35, 2021.
\newblock doi:10.1145/3453721.
\newblock \url{https://arxiv.org/abs/1607.03084}.

\bibitem[Flaxman et~al.(2005)]{FKM05}
Abraham D. Flaxman, Adam Tauman Kalai, and H. Brendan McMahan.
\newblock Online convex optimization in the bandit setting: gradient descent
without a gradient.
\newblock In \emph{Proceedings of the Sixteenth Annual ACM--SIAM Symposium on
Discrete Algorithms}, pages 385--394, 2005.
\newblock \url{https://arxiv.org/abs/cs/0408007}.

\bibitem[Fokkema et~al.(2024)]{Fokkema24}
Hidde Fokkema, Dirk van der Hoeven, Tor Lattimore, and Jack J. Mayo.
\newblock Online Newton method for bandit convex optimisation (extended abstract).
\newblock In \emph{Proceedings of the 37th Conference on Learning Theory},
volume 247 of \emph{Proceedings of Machine Learning Research}, pages 1713--1714,
2024.
\newblock \url{https://proceedings.mlr.press/v247/fokkema24a.html}.
\newblock Full version: arXiv:2406.06506.

\bibitem[Hazan and Levy(2014)]{HL14}
Elad Hazan and Kfir Y. Levy.
\newblock Bandit convex optimization: towards tight bounds.
\newblock In \emph{Advances in Neural Information Processing Systems},
volume 27, pages 784--792, 2014.
\newblock \url{https://proceedings.neurips.cc/paper_files/paper/2014/hash/3ce3c4f8183442957d051dc8ebed2668-Abstract.html}.

\bibitem[Ito(2020)]{Ito20}
Shinji Ito.
\newblock An optimal algorithm for bandit convex optimization with
strongly-convex and smooth loss.
\newblock In \emph{Proceedings of the Twenty Third International Conference on
Artificial Intelligence and Statistics}, volume 108 of \emph{Proceedings of
Machine Learning Research}, pages 2229--2239, 2020.
\newblock \url{https://proceedings.mlr.press/v108/ito20a.html}.

\bibitem[Lattimore(2019)]{Lat20}
Tor Lattimore.
\newblock Improved regret for zeroth-order adversarial bandit convex optimisation.
\newblock \emph{Mathematical Statistics and Learning}, 2(3/4):311--334, 2019.
\newblock doi:10.4171/MSL/17. Published online in 2020.
\newblock \url{https://ems.press/journals/msl/articles/17218}.

\bibitem[Liu et~al.(2025)]{Liu25}
Xiaoqi Liu, Dorian Baudry, Julian Zimmert, Patrick Rebeschini, and Arya Akhavan.
\newblock Non-stationary bandit convex optimization: a comprehensive study.
\newblock In \emph{Advances in Neural Information Processing Systems},
volume 38, 2025.
\newblock doi:10.52202/085713-2198.
\newblock \url{https://arxiv.org/abs/2506.02980}.

\bibitem[Lov\'asz and Vempala(2007)]{LV07}
L\'aszl\'o Lov\'asz and Santosh Vempala.
\newblock The geometry of logconcave functions and sampling algorithms.
\newblock \emph{Random Structures \& Algorithms}, 30(3):307--358, 2007.
\newblock \url{https://faculty.cc.gatech.edu/~vempala/papers/logcon.pdf}.

\bibitem[Rajaraman and Han(2026)]{RH26}
Nived Rajaraman and Yanjun Han.
\newblock The price of hidden curvature: improved lower bounds for bandit convex
optimization.
\newblock arXiv:2607.18652v3, August 20, 2026.
\newblock \url{https://arxiv.org/abs/2607.18652v3}.

\bibitem[Shamir(2013)]{Shamir13}
Ohad Shamir.
\newblock On the complexity of bandit and derivative-free stochastic convex
optimization.
\newblock In \emph{Proceedings of the 26th Annual Conference on Learning Theory},
volume 30 of \emph{Proceedings of Machine Learning Research}, pages 3--24, 2013.
\newblock \url{https://proceedings.mlr.press/v30/Shamir13.html}.

\bibitem[Suggala et~al.(2024)]{Suggala24}
Arun Suggala, Y. Jennifer Sun, Praneeth Netrapalli, and Elad Hazan.
\newblock Second order methods for bandit optimization and control.
\newblock In \emph{Proceedings of the 37th Conference on Learning Theory},
volume 247 of \emph{Proceedings of Machine Learning Research}, pages 4691--4763,
2024.
\newblock \url{https://proceedings.mlr.press/v247/suggala24a.html}.

\bibitem[Yu et~al.(2026)]{Yu26}
Hang Yu, Yu-Hu Yan, and Peng Zhao.
\newblock Improved dimension dependence for bandit convex optimization with
gradient variation.
\newblock In \emph{Proceedings of the 43rd International Conference on Machine
Learning}, volume 306 of \emph{Proceedings of Machine Learning Research},
pages 151442--151483, 2026.
\newblock \url{https://proceedings.mlr.press/v306/yu26bh.html}.
\end{thebibliography}

\clearpage
\appendix
\section{Analytic details and well-definedness}\label{app:analytic}
\subsection{Regularity on the closed domain}\label{app:regularity}
The assumptions \eqref{eq:sc}--\eqref{eq:sm} imply
\begin{equation}\label{eq:gradlip}
 \norm{\nabla f(x)-\nabla f(y)}_2\leq\beta\norm{x-y}_2
 \qquad(x,y\in K).
\end{equation}
To see this using only regularity on $K$, proceed as follows. On
$\operatorname{int}K$, both $f$ and $\beta\norm{x}_2^2/2-f(x)$ are convex.
Convolve locally with a smooth mollifier along a compact interior line segment.
The resulting Hessians are between $0$ and $\beta I_d$, so integrating the
Hessian along the segment gives \eqref{eq:gradlip} for the mollifications.
Pass to the $C^1$ limit, and then approach boundary points from the interior.
The $C^1$ neighborhood assumption ensures continuity of the gradients there.
If $r_0=(4d)^{-1}$, convexity and $0\leq f\leq1$ give, for every unit vector $v$,
\[
 r_0\ip{\nabla f(0)}v\leq f(r_0v)-f(0)\leq1.
\]
Taking a supremum and applying \eqref{eq:gradlip} yields
\begin{equation}\label{eq:globalGappendix}
 \sup_{x\in K}\norm{\nabla f(x)}_2\leq G_0:=4d+\beta.
\end{equation}
This bound is used only for the comparator-distribution error, where its
polynomial dimension dependence enters a logarithm.


\subsection{Covariance contraction on arbitrary convex supports}
\begin{proof}[Covariance part of Lemma~\ref{lem:state}]
The Hessian of $F_t$ is $\lambda tI_d$. To justify the covariance bound without
assuming a smooth boundary, consider for $a>0$ the log-partition function
\[
 \Lambda_a(v)=\log\int_K\exp(-a\norm{x}_2^2/2+\ip vx)\,dx.
\]
Completing the square gives
\[
 \Lambda_a(v)-\frac{\norm v_2^2}{2a}
 =\log\int_K\exp\bigl(-a\norm{x-v/a}_2^2/2\bigr)\,dx.
\]
The right side is concave in $v$ by preservation of log-concavity under
marginalization (Pr\'ekopa's theorem; see~\cite[Theorem 5.1]{LV07}).
All derivatives exist by differentiation of a compactly supported integral,
and $\nabla^2\Lambda_a(v)$ is the covariance under the corresponding tilted
truncated Gaussian. Thus $\nabla^2\Lambda_a\preceq I_d/a$.
Apply this with $a=\eta\lambda t$, then sum $1/t\leq1+\log T$.
The same covariance contraction is used in~\cite[Lemma 2]{Ito20}.
\end{proof}


\subsection{Finite-horizon integrability and measurable kernels}
\begin{proof}[Well-definedness and integrability in Lemma~\ref{lem:state}]
Write $F_t(x)=\lambda t\norm{x}_2^2/2+\ip{b_t}x+c_t$, where
$b_1=0$ and $b_{t+1}=b_t+\widehat g_t-\lambda m_t$.
For a fixed $t$, the covariance of the density proportional to
$\exp(-\eta\lambda t\norm{x}_2^2/2-\eta\ip bx)$ on $K$ is continuous in $b$
and positive definite at every finite $b$. Consequently, on every compact
set of $b$ its minimum eigenvalue has a strictly positive minimum.
Induct on $t$. If $\norm{b_t}_2\leq B_t$ holds uniformly over all histories,
there is a deterministic $\epsilon_t>0$ such that
$\Sigma_t\succeq\epsilon_t I_d$ for all those histories. Then
\eqref{eq:boundedlocal} implies
\[
 \norm{\widehat g_t}_2\leq16s\sqrt{d/\epsilon_t},\qquad
 \norm{b_{t+1}}_2\leq B_t+16s\sqrt{d/\epsilon_t}+\lambda.
\]
This starts from $B_1=0$ and supplies finite deterministic bounds at every
finite $t$. Since $K$ is bounded, the surrogates and additive constants are
also bounded on every finite run. These bounds may be extremely large;
they are used only for integrability, not for a regret or runtime estimate.

The integrals defining the means and covariances are measurable functions
of the finite-dimensional tilt. The gauge is measurable in its parameters,
and one may choose Borel versions of its a.e. partial derivatives by taking
limits of difference quotients and assigning zero where they do not exist.
Thus the displayed densities and estimator are measurable kernels. Since
$q_t$ is absolutely continuous, the exceptional nondifferentiability set
has conditional probability zero at every round. Finite unions remain null.
\end{proof}


\subsection{Proof of the local exponential-moment inequality}
\begin{proof}[Proof of Lemma~\ref{lem:mgf}]
Set $W=\ip v{U-m}$ and $\sigma^2=\E W^2=v^\top\Sigma v$.
If $\sigma=0$ then $W=0$ a.s. Otherwise $W/\sigma$
is a centered one-dimensional log-concave random variable of variance one.
Lov\'asz--Vempala~\cite[Lemma 5.7]{LV07} gives
$\Pp(|W|/\sigma>r)\leq e^{1-r}$ for $r>1$. The same upper bound is valid for
$0\leq r\leq1$ because its right side is then at least one. Hence
\[
 \E|W|^k
 =k\int_0^\infty t^{k-1}\Pp(|W|>t)\,dt
 \leq e\,k!\sigma^k\qquad(k\geq1).
\]
For $\sigma\leq1/2$, the exponential series is absolutely integrable. As
$\E W=0$,
\[
 \E e^W\leq1+\sum_{k=2}^{\infty}\frac{\E|W|^k}{k!}
 \leq1+\frac{e\sigma^2}{1-\sigma}\leq1+2e\sigma^2.
\]
Take logarithms. Log-concavity of linear marginals follows from
\cite[Theorem 5.1]{LV07}.
\end{proof}


\section{Proof of the quadratic lower bound}\label{app:lower}
We use the convention $H^2(P,Q)=\int(\sqrt p-\sqrt q)^2$ for squared Hellinger
distance. Its affinity is $\mathcal A(P,Q)=1-H^2(P,Q)/2$.

\begin{lemma}[Energy-weighted adaptive Hellinger recursion]\label{lem:hellinger}
Consider two experiments $P,Q$ with the same initial private-randomness law
and the same adaptive action rule. Suppose that at every common history $h$
and current action $x$, the squared Hellinger distance between the conditional
observation laws is at most $\varepsilon_t(h,x)\geq0$. Then the laws of the
augmented histories after $T$ observations satisfy
\begin{equation}\label{eq:weightedhellinger}
 H^2(P_{1:T},Q_{1:T})\leq\frac12\sum_{t=1}^T
 \left\{\E_P\varepsilon_t(\mathcal H_{t-1},X_t)
        +\E_Q\varepsilon_t(\mathcal H_{t-1},X_t)\right\}.
\end{equation}
Including the next action does not increase the distance.
\end{lemma}
\begin{proof}
Include the private random seed in the history, so both experiments use a
common action kernel at a common history. If $p,q$ denote the history
densities relative to a common dominating measure, the exact affinity
recursion gives
\[
 H^2(P_{t},Q_{t})-H^2(P_{t-1},Q_{t-1})
 =\int\sqrt{p(h)q(h)}\!\int A_t(dx\mid h)
        H^2(P_t(\cdot\mid h,x),Q_t(\cdot\mid h,x))\,dh.
\]
Bound the inner distance by $\varepsilon_t(h,x)$, then use
$\sqrt{pq}\leq(p+q)/2$. Summing the resulting inequalities proves
\eqref{eq:weightedhellinger}; the initial seed has zero distance. Common
action kernels preserve affinity when the previous history is retained.
No mutual absolute continuity of the conditional laws is needed.
\end{proof}

\Needspace{5\baselineskip}
\begin{proof}[Proof of Theorem~\ref{thm:lowergeneral}]
Use the parameters and random loss sequence in
\eqref{eq:quadscale}--\eqref{eq:lowerfamily}.

\paragraph{Admissibility and exact regret identity.}
For every $x\in RB_2^d$, $|a\ip{\theta}x|\leq aR\sqrt d\leq\Delta/8$.
Thus
\[
 0\leq f_t(x)\leq\frac{\mu R^2}{2}+2aR\sqrt d+2\sigma
 \leq1-\Delta+\frac\Delta4+\frac\Delta4
 =1-\frac\Delta2\leq1.
\]
Every realized loss is smooth on all of $\R^d$ with Hessian exactly $\mu I_d$.
Moreover,
$\norm{x_\theta^*}_2^2=da^2/\mu^2\leq R^2/64$.
The common minimizer is therefore feasible and strictly interior. Pointwise,
\begin{equation}\label{eq:quadregret}
 f_t(x)-f_t(x_\theta^*)=\frac\mu2\norm{x-x_\theta^*}_2^2.
\end{equation}
All offsets are presampled independently of the learner's randomness, and
observations are exact values of these selected loss functions.

\paragraph{One-step distinguishability.}
Let $\theta^{(i)}$ flip sign $i$. The difference between the observation
locations at a common $x$ is $2a x_i$ in absolute value. The square root
$\sqrt{q_0(z)}=\cos(\pi z/2)\mathbf1_{\{|z|\leq1\}}$, extended by
zero, belongs to the Sobolev space $H^1(\R)$ of square-integrable functions
with square-integrable weak derivative. Its
squared derivative norm is $\pi^2/4$. The translation inequality for $H^1$
functions, obtained by integrating the weak derivative and applying
Cauchy--Schwarz, gives
\[
 H^2(q_0(\cdot-u),q_0(\cdot-v))\leq\frac{\pi^2}{4}(u-v)^2.
\]
After scaling the offset by $\sigma$, the one-step squared distance is at most
$\pi^2a^2x_i^2/\sigma^2$. If $P_\theta$ is the full augmented-history law,
Lemma~\ref{lem:hellinger} yields
\begin{equation}\label{eq:pairenergy}
 H^2(P_\theta,P_{\theta^{(i)}})
 \leq\frac{\pi^2a^2}{2\sigma^2}\sum_{t=1}^T
   \left(\E_\theta X_{t,i}^2+\E_{\theta^{(i)}}X_{t,i}^2\right).
\end{equation}

\paragraph{Regret controls query energy.}
Let $\E$ also average over the uniform prior on $\theta$, and put
\[
 \mathsf R=\E\Reg_T(x_\theta^*),\qquad
 \mathsf D=\frac1d\sum_{i=1}^d2^{-d}\sum_\theta
                    H^2(P_\theta,P_{\theta^{(i)}}).
\]
Averaging \eqref{eq:pairenergy}, and using the sign-flip bijection, gives
\[
 \mathsf D\leq\frac{\pi^2a^2}{d\sigma^2}
                     \sum_{t=1}^T\E\norm{X_t}_2^2.
\]
Since $\norm{X_t}_2^2\leq2\norm{X_t-x_\theta^*}_2^2
+2\norm{x_\theta^*}_2^2$, \eqref{eq:quadregret} implies
\begin{equation}\label{eq:energyregret}
 \mathsf D\leq\frac{\pi^2a^2}{d\sigma^2}
       \left(\frac{4\mathsf R}{\mu}
                          +\frac{2dTa^2}{\mu^2}\right).
\end{equation}

\paragraph{Testing the signs.}
For each coordinate, condition on all other signs and pair the two full
history laws. The uniform-prior binary error is at least $(1-\TV)/2$.
Although the action at time $t$ uses less than the full history, the same
lower bound holds for a sign prediction extracted from that action, by data
processing. For example, predict $+1$ when $X_{t,i}\geq0$ and $-1$ otherwise.
Its error event is contained in $\{\theta_iX_{t,i}\leq0\}$, so ties
cause no difficulty. As $\TV(P,Q)\leq H(P,Q)$, Cauchy--Schwarz gives
\begin{equation}\label{eq:signsum}
 \sum_{t=1}^T\sum_{i=1}^d\Pp(\theta_iX_{t,i}\leq0)
 \geq\frac{dT}{2}\bigl(1-\sqrt{\mathsf D}\bigr).
\end{equation}
Each wrong-sign coordinate contributes at least $a^2/(2\mu)$ in
\eqref{eq:quadregret}. Hence
\begin{equation}\label{eq:testingregret}
 \mathsf R\geq\frac{dTa^2}{4\mu}\bigl(1-\sqrt{\mathsf D}\bigr).
\end{equation}

\paragraph{Closing the tradeoff.}
Write $\mathsf B=dTa^2/\mu$. If $\mathsf R<\mathsf B/8$, then
\eqref{eq:energyregret} and \eqref{eq:quadscale} imply
\[
 \mathsf D<\frac{5\pi^2Ta^4}{2\mu^2\sigma^2}
 \leq\frac5{512}<\frac14.
\]
Equation~\eqref{eq:testingregret} would then give
$\mathsf R\geq\mathsf B/8$, a contradiction. Therefore
\begin{align*}
 \mathsf R\geq\frac{dTa^2}{8\mu}
 &=\min\!\left\{\frac{MT}{512},\;
                 \frac{\Delta^2T}{512M},\;
                 \frac{\Delta d\sqrt T}{1024\pi}\right\}\\
 &\geq\frac1{1024\pi}\min\{QT,\Delta d\sqrt T\}.
\end{align*}
Finally, the supremum over admissible deterministic preselected sequences is
at least this average over $\theta,\xi_{1:T}$ for any fixed randomized learner.
This proves the theorem under the same exact one-point feedback model.
\end{proof}

\section{Uniformity at fixed curvature}\label{app:fixed}
\begin{proof}[Proof of Corollary~\ref{cor:fixed}]
First suppose $0<\mu\leq1$ and use $R=1$. Then $M=\mu$,
$\Delta=1-\mu/2\geq1/2$, and
$Q=\min\{\mu,\Delta^2/\mu\}\geq\mu/4$. It follows that
\[
 \min\{QT,\Delta d\sqrt T\}
 \geq\frac14\min\{\mu T,d\sqrt T\},
\]
which gives \eqref{eq:unitlower}. At fixed $\mu>0$ this implies a positive
constant multiple of $\min\{T,d\sqrt T\}$.

For completeness, arbitrary fixed $\mu>0$ can be treated uniformly over all
integer dimensions satisfying $\mu<32d^2$. Define the integer
\[
 D_\mu=\lfloor\sqrt{\mu/32}\rfloor+1
\]
and the fixed radius
\[
 R_\mu=\max\!\left\{\frac1{4D_\mu},\;
                               \min\{1,\mu^{-1/2}\}\right\}.
\]
For every such $d$, $d\geq D_\mu$ and
$R_\mu\in[(4d)^{-1},1]$. Moreover,
$M_\mu=\mu R_\mu^2<2$: the first radius term gives
$\mu/(16D_\mu^2)<2$, and the second gives $\min\{\mu,1\}\leq1$.
Let $\Delta_\mu=1-M_\mu/2$ and
$Q_\mu=\min\{M_\mu,\Delta_\mu^2/M_\mu\}$. The constant
\[
 c_\mu=\frac1{1024\pi}\min\{Q_\mu,\Delta_\mu\}>0
\]
depends on $\mu$ alone. Theorem~\ref{thm:lowergeneral} on $R_\mu B_2^d$ then
proves the lower inequality of \eqref{eq:fixedmatching}.

For fixed $\kappa=\beta/\mu$, the bracket in \eqref{eq:generalupper} is at
most $C_\kappa\log^2(ed)\log(edT)$. This proves the upper inequality.
\end{proof}

\paragraph{What is and is not matched.}
The fixed-pair statement allows constants depending on $\mu,\beta$, but not
on $d,T$; only logarithms of $d,T$ are suppressed. It is not obtained by
silently scaling all losses: such a scaling would also change their range and
the signal-to-offset ratio. The general lower construction instead retains
curvature, radius, and bounded-loss headroom simultaneously. It stays
nontrivial when $\beta=\mu$, but does not establish a growing $\kappa$ lower
factor. Joint optimality in $(d,T,\mu,\beta)$ remains distinct from
\eqref{eq:fixedmatching}. The zero-regret endpoint of
Proposition~\ref{prop:admissibility} also prevents a uniform positive lower
constant over all feasible curvature parameters.

\end{document}
